Sanguosuo demonstrates that mobile security is an end-to-end property. A polished Android form cannot compensate for an API that retains connections without limits, combines untrusted text with SQL, stores recoverable passwords, or sends credentials over an unencrypted channel. The appropriate controls operate at different boundaries: timeouts, connection limits, rate limits, and upstream filtering preserve availability; parameterized queries and least privilege preserve database integrity; and adaptive salted hashing combined with TLS protects credentials.

The prototype already applies two central defenses consistently. Request values are bound through SQL placeholders, and account passwords are stored as bcrypt hashes with a cost of 10 and checked through the bcrypt library. The database schema reinforces these choices with constraints. Conversely, the repository does not yet implement or measure explicit Slowloris mitigation, and several integration defects weaken its authentication story: the Android client uses cleartext HTTP, issued session tokens are incompatible with JWT validation, a hard-coded bypass exists, registration omits an \verb|await|, and automated tests do not cover the security behavior.

The next milestone is therefore concrete: secure the release transport with HTTPS, choose and complete one session model, remove test bypasses and detailed production errors, await database writes, configure and measure HTTP resource limits, add endpoint-specific throttling, and build automated security regression tests. Once those tasks and the controlled experiments in Section \ref{demo} are complete, the project can support empirical conclusions about resilience rather than relying on design intent alone.
